% Image credits for the photographs and AI illustrations of Book 3
% (University Chemistry, Year 2). Machine-readable license records live in
% images/book3/CREDITS.md; keep the two files in sync. The Book 3 writer
% owns this file: add one \item per photograph (in an itemize, after the
% first paragraph) as photographs land.
\chapter*{Image Credits}

The drawings in this book are original. Photographs, where used, are used
under their respective free licenses, with thanks to their authors; full
source links and license URLs for every photograph are listed in
\texttt{images/book3/CREDITS.md} in the book's repository.
The hazard pictograms are those of the Globally Harmonized System of
Classification and Labelling of Chemicals, as drawn for the
\texttt{ghsystem} package of \TeX\ Live.

\bigskip
Alongside the photographs and the original drawings, some figures are
AI-generated illustrations, produced with OpenAI image generation from
prompts written by the book's editors, and reviewed for chemical
accuracy before inclusion. The full prompt for every generated image is
recorded in \texttt{images/book3/ai/PROMPTS.md} in the repository.

\bigskip
\noindent Photographs, by chapter:
\begin{itemize}\small
  \item Ch.~1: Germain Henri Hess, portrait by I.~I. Reimers (1837--1838),
        CC BY-SA 4.0; Marcellin Berthelot, unknown author, public domain
        (both Wikimedia Commons).
  \item Ch.~2: Josiah Willard Gibbs, unknown photographer, public domain
        (Wikimedia Commons).
  \item Ch.~3: François-Marie Raoult, unknown photographer, public domain
        (Wikimedia Commons).
  \item Ch.~4: Jacobus Henricus van 't Hoff, by Nicola Perscheid (1904),
        public domain; Henry Le Chatelier, unknown author, CC BY-SA 4.0;
        Fritz Haber and Carl Bosch, Nobel Foundation, public domain (all
        Wikimedia Commons).
  \item Ch.~6: blast furnace at Duisburg-Nord, by Dietmar Rabich, CC BY-SA
        4.0; thermite welding of a rail, by PetrS., CC BY-SA 3.0 (both
        Wikimedia Commons).
  \item Ch.~7: distillation column of a refinery, by CEphoto, Uwe Aranas,
        CC BY-SA 3.0 (Wikimedia Commons).
  \item Ch.~9: Michael Faraday, by Thomas Phillips (1842), public domain
        (Wikimedia Commons).
  \item Ch.~10: Jaroslav Heyrovsk\'y, archive of the J.~Heyrovsk\'y Institute
        of Physical Chemistry, CC BY-SA 3.0 (Wikimedia Commons).
  \item Ch.~11: a voltaic pile, by GuidoB, CC BY-SA 3.0; Charles Martin Hall
        and Paul H\'eroult, unknown authors, public domain (all Wikimedia
        Commons).
  \item Ch.~12: a copper statue, by Elcobbola, public domain; a sacrificial
        anode, by Zwergelstern, CC BY-SA 3.0 (both Wikimedia Commons).
  \item Ch.~13: Erwin Schr\"odinger, Nobel Foundation (1933), public domain
        (Wikimedia Commons).
  \item Ch.~14: liquid oxygen in a magnet, by Pieter Kuiper, public domain;
        Robert Mulliken and Friedrich Hund, by GFHund, CC BY 3.0 (both Wikimedia
        Commons).
  \item Ch.~16: August Kekul\'e (1873), unknown author, public domain
        (Wikimedia Commons).
  \item Ch.~18: transition-metal solutions, by Benjah-bmm27, public domain; Alfred
        Werner, Les Prix Nobel (1914), public domain (both Wikimedia Commons).
  \item Ch.~19: a ruby crystal, public domain; an emerald, by Didier Descouens,
        CC BY-SA 3.0; copper sulfate crystals, by Stephanb, CC BY-SA 3.0 (all
        Wikimedia Commons).
  \item Ch.~20: Akira Suzuki, Ei-ichi Negishi and Richard Heck, by BloodIce, CC BY-SA
        4.0 (Wikimedia Commons).
  \item Ch.~22: Charles Friedel (1890s) and James Mason Crafts (\textit{Popular
        Science Monthly}, 1900), public domain (Wikimedia Commons).
  \item Ch.~24: Michel Eug\`ene Chevreul, by Nicolas Eustache Maurin, public domain
        (Wikimedia Commons).
  \item Ch.~25: Alexander Borodin, unknown author, public domain; Ludwig Claisen, by
        Veronica.villagrasa, CC BY-SA 4.0 (both Wikimedia Commons).
  \item Ch.~26: Robert Robinson, unknown author, public domain (Wikimedia Commons).
  \item Ch.~28: Elias James Corey, by Trvthchem, CC0 (Wikimedia Commons).
  \item Ch.~29: Wallace Carothers in the laboratory, unknown photographer, public domain
        (Wikimedia Commons).
  \item Ch.~30: Emil Fischer, Atelier Victoria, Berlin, around 1895, public domain (Wikimedia
        Commons).
  \item Ch.~31: Mikhail Tsvet, unknown author, public domain (Wikimedia Commons).
  \item Ch.~32: Francis William Aston (1922), unknown author, public domain; Gustav Kirchhoff,
        Robert Bunsen and Henry Roscoe (1862), public domain (both Wikimedia Commons).
  \item Ch.~34: William Sealy Gosset (1908), public domain (Wikimedia Commons).
%% next chapter here
\end{itemize}
\clearpage
